\section{社会应对：白领通缩之后怎么办}

本期最沉重的部分，是白领通缩和失业窗口。技术上，Coding 模型压缩研发周期；组织上，少数人能完成更多任务；经济上，白领劳动价格会被重新定价。问题不是“AI 会不会替代人”，而是社会能否给被替代或被重组的人提供新位置。

\subsection{三类人会先受到冲击}

\noindent\begin{tabularx}{\textwidth}{p{0.22\textwidth}p{0.34\textwidth}X}
\toprule
人群 & 为什么先受影响 & 应对方向 \\
\midrule
重复性知识工作者 & 工作可被文本、表格、代码表达 & 转向任务定义、审核、客户理解和流程设计。 \\
初级程序员 & 大量实现代码可由模型生成 & 学会架构、测试、review、系统理解。 \\
中层协调者 & 信息转发和进度管理可被工具自动化 & 转向判断、资源配置、跨团队冲突解决。 \\
\bottomrule
\end{tabularx}

\begin{knowledgebox}{再培训的真正目标}
再培训不应只是教人“用 AI 工具”，而应训练人定义问题、验证结果、组织工作流和理解业务。工具会变，判断力和责任结构更慢变。
\end{knowledgebox}

\subsection{本章小结}

白领通缩是技术扩散速度超过组织吸收速度的结果。社会应对不能只靠个人学习，也需要企业组织和制度设计重新分配生产力红利。


